\bookchapter{Capítulo 8: Comunicação: mensagens, chamadas e redes sociais}

\chapterintro{Comunicar-se é a missão mais antiga e central do celular. Hoje, isso já evoluiu de chamadas e SMS simples para um ecossistema complexo que inclui mensagens instantâneas, redes sociais e e-mail.}

\booksection{Recursos avançados do app de telefone}

O aplicativo de telefone dos celulares modernos já vai muito além de um simples discador: integra ferramentas poderosas como gravação de chamadas, bloqueio de spam e identificação inteligente.

\booksubsection{Gravação de chamadas (atenção à conformidade legal):}

\begin{itemize}
  \item Para que serve: registrar o conteúdo de chamadas importantes, usado em atas de reunião, entrevistas, promessas de clientes ou como prova em disputas.
  \item Aviso legal: na China, usar a gravação de chamadas como prova exige o consentimento da outra parte, ou que seja usada para defender os seus próprios direitos legítimos.
  \item Gravação unilateral vs. consentimento mútuo: no plano jurídico, uma gravação feita "sem o consentimento da outra parte" não é automaticamente inválida. Nas interpretações judiciais relevantes do Supremo Tribunal Popular, desde que a gravação não viole gravemente os direitos legítimos de terceiros, não infrinja proibições legais nem ofenda a moral e os bons costumes, ainda pode ser apresentada como prova. Porém, tribunais de algumas províncias adotam critérios mais rigorosos para "gravações secretas" em processos cíveis; em casos criminais, a revisão de legalidade costuma ser ainda mais estrita.
  \item Diferenças regionais: tribunais de Pequim, Xangai, Guangdong e outros lugares são mais flexíveis com gravações unilaterais como prova cível; já gravações que envolvam segredos comerciais ou privacidade pessoal exigem muita cautela em qualquer província, e podem ser excluídas por violação da privacidade.
  \item Como fazer: no Android, a tela de chamada costuma ter um botão "Gravar" bem visível; o iOS não oferece gravação nativa de chamadas, então é preciso usar um app de terceiros ou um fone Bluetooth com função de gravação.
\end{itemize}

\begin{tcbwarn}{Limites no uso de gravações}
\begin{itemize}
  \item Não use gravações para chantagem, divulgação ou humilhação pública de outras pessoas.
  \item Conteúdo que envolva a privacidade de terceiros (por exemplo, uma conversa alheia gravada por acaso) não pode ser espalhado livremente.
  \item Em situações importantes, recomenda-se avisar verbalmente no início da chamada que "esta conversa será gravada" — isso é mais correto e dá mais força probatória.
\end{itemize}

\end{tcbwarn}

\booksubsection{Conteúdo aprofundado sobre combate a chamadas indesejadas:}

\begin{itemize}
  \item Como funciona: o app de telefone consulta uma enorme base de números pela internet e, na hora da chamada, mostra "spam", "golpe" ou o nome da empresa.
  \item Como ativar: no Android, abra o app de telefone > Configurações > Bloqueio de spam; no iOS, vá em Configurações > Telefone > Bloqueio e Identificação de Chamadas. Os caminhos detalhados por marca e as opções recomendadas estão na tabela comparativa do Capítulo 13 «Ferramentas úteis do sistema».
  \item Principais plataformas de marcação: na China, a marcação de números é feita principalmente por três provedores de dados — Tencent Mobile Manager, 360 Mobile Guard e Diànghuà Bāng (Tapir). As operadoras (China Mobile/Unicom/Telecom) também costumam usar esses dados para exibir o "nome do titular" e o bloqueio de spam. As bases de marcação dos fabricantes (Huawei, Xiaomi, OPPO, vivo) geralmente se integram a esses provedores.
  \item E se o seu número for marcado por engano: se o seu número de trabalho normal for marcado erradamente como "delivery", "correios" ou "spam", você pode recorrer pelos seguintes canais: | | Diànghuà Bāng: acesse a plataforma de recurso de números (haoma.tk.cn) e envie um documento de identidade.
  \item 360: envie pelo recurso "Recurso de número" dentro do app 360 Mobile Guard.
  \item Tencent: envie pela página "Recurso de número" do miniapp "Tencent Mobile Manager".
  \item Três grandes operadoras: usuários da China Mobile podem enviar o SMS "KTFSR" para 10086 e ativar o antispam; usuários da Telecom devem seguir a conta oficial "Tianyi Antispam"; usuários da Unicom configuram no "Unicom Mobile Manager".
\end{itemize}

Antispam avançado: ative o "bloqueio de chamadas de alta frequência" da operadora (por exemplo, a limitação de chamadas de alta frequência da China Mobile) para bloquear automaticamente chamadas de voz por IA.

Recursos comuns de bloqueio nativo do celular: | | Identificação de chamadas / bloqueio inteligente: marca e bloqueia automaticamente spam, golpes e telemarketing com base em bases de números.

Lista negra / lista branca: bloqueie manualmente números de spam conhecidos; adicione contatos e números de entrega à lista branca para evitar bloqueios acidentais.

Filtro de palavras-chave em SMS: mensagens com termos como "prêmio", "empréstimo" ou "clique no link" vão automaticamente para a caixa de bloqueio.

Silenciar chamadas desconhecidas no iOS: chamadas fora dos contatos são apenas silenciadas e vão para o correio de voz, sem interromper você; ideal para quem recebe muito spam, mas fique atento a chamadas importantes de desconhecidos.

Integração com o modo Não Perturbe: à noite ou em reuniões, permita apenas chamadas dos contatos, combinado com agendamento automático — a qualidade do sono melhora na hora.

Como combinar com apps úteis: priorize o bloqueio nativo do sistema e some o Centro Nacional Antifraude e o antispam da operadora; para limpeza, ferramentas leves como Qingzhuo são opções, mas instale apenas um app de segurança/limpeza de cada tipo — veja o Capítulo 13.

\booksubsection{O que é VoLTE/VoNR e como ativar:}

\begin{itemize}
  \item VoLTE (Voice over LTE): tecnologia de chamadas de voz sobre a rede 4G; durante a chamada o aparelho não cai para 2G/3G, permite falar e usar a internet ao mesmo tempo, com qualidade de áudio melhor e conexão mais rápida.
  \item VoNR (Voice over NR): chamadas de voz sobre 5G, a próxima evolução do VoLTE. Latência menor, qualidade de áudio melhor e suporte a videochamadas em alta definição. Exige três coisas: celular 5G + plano 5G + cobertura 5G local.
  \item Como ativar: no iOS, em "Configurações > Dados Móveis > Opções de Dados > Voz e Dados", ative "VoLTE"; no Android, em "Configurações > Rede móvel > Dados móveis > Avançado", procure a opção "Chamada HD VoLTE". As operadoras costumam ativar por padrão; se a internet cair durante a chamada, provavelmente o VoLTE não está ativado.
\end{itemize}

\begin{tcbinfo}{Curiosidade}
Se a internet do celular ficar lenta ou cair durante uma chamada, é quase certo que o VoLTE/VoNR não está ativado e o aparelho caiu para 2G/3G. Ative para “falar no telefone e ver vídeos ao mesmo tempo”.

\end{tcbinfo}

\booksection{SMS e iMessage/RCS}

O SMS tradicional não foi aposentado: continua sendo um canal confiável para receber códigos de verificação, avisos de bancos e mensagens oficiais.

\begin{itemize}
  \item iMessage: mensagens enviadas pela internet entre dispositivos Apple, exibidas em balões azuis. Suporta fotos e vídeos de alta qualidade, confirmação de leitura etc., é totalmente gratuito e com criptografia de ponta a ponta.
  \item RCS: protocolo de comunicação de nova geração impulsionado por operadoras e pelo Google em todo o mundo. Funciona de forma semelhante ao iMessage ou WeChat, mas baseado no número de telefone.
  \item Interoperabilidade entre iMessage e RCS: a partir do iOS 18, a Apple passou a suportar o padrão RCS. Isso significa que iPhones e Androids podem trocar imagens em alta definição, confirmações de leitura e indicadores de digitação via RCS — fim da experiência dos "balões verdes". O pré-requisito é que as operadoras de ambos os lados tenham RCS habilitado e que o sistema esteja atualizado o suficiente.
\end{itemize}

\booksubsection{Situação do RCS na China:}

\begin{itemize}
  \item China Mobile: implementou o RCS mais cedo; algumas províncias já suportam o serviço "5G Messaging", com mensagens multimídia e cartões interativos. Exige aparelho compatível e ativação na loja da operadora.
  \item China Telecom: oferece o serviço via "Telecom 5G Messaging"; alguns planos já incluem por padrão, e usuários antigos precisam enviar um comando SMS específico para ativar.
  \item China Unicom: implementou 5G Messaging nas principais cidades, mas a cobertura de aparelhos é irregular; muitos Androids ainda exigem ativar manualmente "Mensagens aprimoradas" ou "5G Messaging" no app de SMS.
  \item Experiência real: a adoção do RCS na China foi mais lenta que o previsto, principalmente pela força do ecossistema WeChat e pela falta de padronização nos aparelhos. No dia a dia, o WeChat continua dominando; o RCS assume mais o papel de "canal oficial" (avisos de bancos, mensagens governamentais).
\end{itemize}

\begin{tcbwarn}{Limites de segurança com códigos de verificação}
\begin{itemize}
  \item Nunca encaminhe códigos de verificação: qualquer pessoa que peça para você "ler o código" ou "encaminhar o SMS" pode ser golpista — inclusive quem se diz "atendente", "entregador" ou "polícia/justiça". Código de verificação = o seu dinheiro.
  \item Cuidado com SMS de phishing: desconfie de links em mensagens como "seu ETC expirou", "cartão de saúde congelado" ou "indenização por encomenda perdida". Instituições sérias nunca pedem número de cartão ou senha por link em SMS.
  \item Identificação do número: SMS oficiais de bancos costumam vir de códigos curtos que começam com 95 ou de 1069 + assinatura da empresa. Qualquer "aviso de banco" vindo de um celular comum de 11 dígitos é quase sempre golpe.
  \item Backup de autenticação em dois fatores: não troque de número sem necessidade se ele estiver vinculado a contas importantes; se precisar trocar, desvincule antes todas as contas críticas. Caso contrário, o novo dono do número pode redefinir suas senhas por SMS.
\end{itemize}

\end{tcbwarn}

\booksection{WeChat/QQ: gerenciando os superapps sociais}

WeChat e QQ já ultrapassaram o papel de ferramentas de comunicação e se tornaram “superapps” que reúnem social, pagamentos e informação.

\booksubsection{Backup e migração do histórico de conversas:}

\begin{enumerate}
  \item Entre no WeChat para PC ou Mac no computador.
  \item Clique no menu do canto inferior esquerdo > [Migrar e fazer backup] > [Backup e restauração].
  \item Escolha "Fazer backup do histórico no computador" e selecione no celular as conversas que deseja salvar.
  \item Migrar para um celular novo: no WeChat do aparelho antigo, vá em [Eu] > [Configurações] > [Geral] > [Migração e backup do histórico] > [Migrar] > [Migrar para celular/tablet], e escaneie o código QR para concluir.
\end{enumerate}

\booksubsection{Métodos avançados de limpeza de espaço:}

\begin{itemize}
  \item Entrada padrão: WeChat [Eu] > [Configurações] > [Geral] > [Armazenamento]. Lá você pode limpar o cache e fazer limpeza fina do histórico por período e tipo de arquivo.
  \item Limpeza dentro da conversa: abra uma conversa > "..." no canto superior direito > [Histórico de conversas] > [Gerenciar armazenamento]. É possível apagar só imagens, vídeos e arquivos de uma conversa, mantendo o texto.
  \item Limpeza de imagens/vídeos originais: em grupos, as mídias em alta definição ocupam mais espaço. Em [Gerenciar armazenamento], marque "Originais" para excluir em lote.
  \item Desativar download automático: [Configurações] > [Geral] > [Fotos, vídeos, arquivos e chamadas], desative "Download automático" para evitar que o Wi-Fi baixe silenciosamente muitos arquivos de grupo.
  \item Cache de miniapps: em [Descobrir] > [Miniapps], pressione e segure miniapps pouco usados para excluí-los e liberar o cache local.
  \item Backup compactado do histórico: ao fazer backup no WeChat para computador, marque "Compactar". Vários GB de histórico podem ser reduzidos a 30\%–50\% do tamanho original.
\end{itemize}

\booksubsection{Dicas de gerenciamento de grupos:}

\begin{itemize}
  \item Anúncio do grupo: o administrador publica em "..." > [Anúncio do grupo], no canto superior direito. Todos recebem alerta forte — ideal para avisos importantes. Aparece automaticamente para novos membros.
  \item Tarefas do grupo: o administrador pode fixar uma mensagem como "tarefa do grupo"; todos veem um aviso vermelho que desaparece ao concluir. Ótimo para inscrições e coleta de tarefas.
  \item Lista colaborativa: digite "\#lista" no campo de texto para iniciar uma lista em sequência, útil para pedidos de comida e contagem de inscrições. O resultado é consolidado no topo do grupo, sem inundar o chat.
  \item Live no grupo: " + " no canto superior direito > [Live do grupo]. A transmissão fica restrita ao grupo — bom para treinamentos internos, encontros de família e testes de produto.
  \item Dobrar o grupo: na lista de grupos, deslize o nome para a esquerda > [Dobrar]. Grupos menos importantes ficam em "Grupos dobrados" e só avisam quando mencionarem você.
  \item Somente administradores @todos: em grupos com mais de 500 membros, membros comuns não podem mencionar todos, evitando distrações.
\end{itemize}

\booksubsection{Gerenciamento fino de permissões do Moments:}

\begin{itemize}
  \item Visibilidade por grupo: ao publicar, toque em "Quem pode ver", escolha "Parcialmente visível" ou "Ocultar de", e use etiquetas previamente criadas ("Colegas", "Família", "Clientes") para segmentar.
  \item Período de visualização dos amigos: [Eu] > [Configurações] > [Permissões de amigos] > [Moments] > [Período que os amigos podem ver]. Opções: "Últimos 3 dias / 1 mês / 6 meses / Tudo". Recomenda-se "6 meses" para equilibrar privacidade e vida social.
  \item Ocultar um post do Moments: pressione e segure a publicação > [Não ver os Moments dele/dela] ou [Deixar de seguir]. A pessoa não fica sabendo.
  \item Gerenciar anúncios do Moments: [Eu] > [Configurações] > [Sobre o WeChat] > [Informações pessoais e permissões] > [Gerenciamento de anúncios personalizados], e desative a recomendação personalizada baseada em comportamento.
  \item Dez posts para estranhos: [Permissões de amigos] > [Moments] > [Permitir que estranhos vejam 10 Moments]. Desativar evita que desconhecidos vejam suas publicações antigas.
\end{itemize}

\booksection{Comparação de videochamadas}

Cada ferramenta de videochamada tem seus pontos fortes. Escolher a certa para cada situação deixa a comunicação mais fluida.

\begin{itemize}
  \item WeChat Video: cobertura nacional, fácil para todas as idades. Ideal para contato diário com família e amigos. Suporta até 20 pessoas (via miniapp "Video em grupo"); a desvantagem é a queda acentuada de qualidade em redes ruins.
  \item FaceTime: exclusivo do ecossistema Apple, com ótima qualidade e criptografia de ponta a ponta. Desde o iOS 15, aceita convites por "link" para usuários Android/Web, mas o anfitrião precisa ser um dispositivo Apple. Ideal para famílias com aparelhos mistos.
  \item Zoom: primeira opção para reuniões internacionais, com reuniões gratuitas de até 40 minutos para 100 pessoas. Tem compartilhamento de tela, fundo virtual e salas de breakout — ótimo para colaboração remota e internacional.
  \item Tencent Meeting: preferida no ambiente corporativo chinês, integra-se profundamente ao WeChat e permite entrada por QR code. Suporta reuniões de 300 pessoas (plano pago), com gravação na nuvem, legendas ao vivo e atas automáticas — ideal para treinamentos e eventos grandes.
  \item Recomendações por cenário: 1 a 1 no dia a dia, WeChat; entre usuários Apple, FaceTime; reuniões de negócios formais, Zoom ou Tencent Meeting (na China, Tencent; no exterior, Zoom).
\end{itemize}

\booksubsection{Referência de largura de banda e consumo de dados:}

\begin{itemize}
  \item Videochamada no WeChat: cerca de 500 Kbps–1 Mbps; qualidade padrão consome 3–5 MB por minuto, HD cerca de 8–12 MB.
  \item FaceTime: SD cerca de 1 Mbps, 720p cerca de 2 Mbps, 1080p exige mais de 4 Mbps. Consumo de 5–15 MB por minuto.
  \item Zoom/Tencent Meeting: 720p cerca de 1,5–3 Mbps, 1080p exige 3–5 Mbps. Cerca de 300–800 MB por hora; o anfitrião de reuniões multitudinárias consome mais.
  \item Dica: em videochamadas longas, prefira Wi-Fi; em rede móvel, reduza a qualidade (desative "Vídeo HD" nas configurações do app) para economizar dados.
\end{itemize}

\booksubsection{Dicas para melhorar a imagem nas videochamadas:}

\begin{itemize}
  \item Iluminação: fique de costas para a janela e de frente para a luz (luz natural é ideal), evitando contraluz que escurece o rosto.
  \item Ângulo: posicione a câmera um pouco acima dos olhos para evitar o queixo duplo; em chamadas na horizontal, segure o celular com as duas mãos.
  \item Áudio: use fones para evitar eco e microfonia; em ambientes silenciosos, não ative o viva-voz.
  \item Rede: priorize Wi-Fi; em 4G/5G, aproxime-se da janela para melhorar o sinal; em reuniões com muita gente, desative as outras câmeras e mantenha só quem fala.
  \item Desfoque de fundo: Tencent Meeting e Zoom oferecem desfoque por IA, útil para esconder fundos bagunçados. Prefira roupas de cor sólida e evite cores parecidas com o fundo, para não parecer "chroma key".
\end{itemize}

\begin{tcbinfo}{Dica rápida}
Se a rede travar, desligue o vídeo e mantenha só o áudio — é mais educado do que desligar a chamada. Antes de reuniões importantes, use o “Teste de rede” do Tencent Meeting para medir a largura de banda.

\end{tcbinfo}

\booksection{E-mail: configurar e usar no celular}

\begin{itemize}
  \item Adicionar conta de e-mail: no iOS, Configurações > Mail > Contas > Adicionar Conta; no Android, costuma haver um app "E-mail".
  \item Dicas: configure notificações diferentes para cada conta; use gestos de deslizar para arquivar, apagar ou marcar e-mails rapidamente.
\end{itemize}

\booksubsection{Diferenças entre IMAP, POP3 e Exchange:}

\begin{itemize}
  \item POP3: baixa os e-mails do servidor e, por padrão, apaga o original, deixando apenas uma cópia local. Bom para quem usa um só dispositivo e precisa ler offline. Desvantagem: ao trocar de aparelho, o histórico some. Já não é recomendado.
  \item IMAP: os e-mails ficam sempre no servidor; o local é só cache. A sincronização entre celular/computador/tablet é consistente — lido em um, lido em todos. É o protocolo mais recomendado hoje.
  \item Exchange (inclui ActiveSync): protocolo da Microsoft que, além de e-mail, sincroniza calendário, contatos e tarefas. Primeira opção para e-mail corporativo. Suportado por Outlook, NetEase Enterprise Mail e Tencent Enterprise Mail.
  \item Recomendação: e-mail pessoal use IMAP; e-mail corporativo, prefira Exchange e, em seguida, IMAP; evite POP3, a menos que precise de arquivamento especial.
\end{itemize}

\booksubsection{Como configurar e-mail corporativo:}

\begin{enumerate}
  \item iOS: [Configurações] > [Mail] > [Contas] > [Adicionar Conta] > escolha "Microsoft Exchange" ou "Outra".
  \item Android: no app "E-mail" nativo ou no Outlook, escolha "Adicionar conta" e informe o endereço completo e a senha.
  \item Preenchimento do servidor (exemplo Exchange): o endereço costuma ser "mail.dominiodaempresa.com" ou "outlook.office365.com"; no campo de domínio, use o nome curto da empresa.
  \item Se aparecer "não foi possível verificar o servidor", preencha manualmente: servidor de entrada (IMAP) porta 993, criptografia SSL; servidor de saída (SMTP) porta 465 ou 587, criptografia TLS.
  \item Após entrar, ative "Push" em [Configurações da conta] para receber e-mails em tempo real, e desative "Pré-visualização de e-mail" para proteger segredos comerciais.
\end{enumerate}

\begin{tcbinfo}{Observação}
E-mails corporativos costumam exigir que o administrador habilite uma “senha de autorização do cliente” ou “código de autorização IMAP/SMTP” — não dá para usar direto a senha do webmail. Se o login falhar, gere o código de autorização no webmail e insira no celular.

\end{tcbinfo}

\begin{tcbpractice}{Prática: fazer backup do histórico do WeChat para o computador}
\begin{enumerate}
  \item Inicie o backup no computador: abra o WeChat para PC > menu no canto inferior esquerdo > [Migrar e fazer backup] > [Backup e restauração].
  \item Escolha backup para o computador: clique em "Fazer backup do histórico no computador".
  \item No celular, selecione o conteúdo do backup e aguarde a conclusão.
\end{enumerate}

\end{tcbpractice}
